%% ma: L12-B05-0014
%% lop: 12
%% chuong: 1
%% bai: 5
%% dang: 12.5.2
%% loai: TN4
%% muc: VD
%% dap_an: "B"
%% nguon: "Ảnh giáo viên gửi 11/10/2026 (Đề kiểm tra giữa HK1 Toán 12, THPT Việt Đức 2026-2027, mã đề 0103), Phần I Câu 4"
%% kien_thuc: "Bài 5 (vận tốc là đạo hàm của quãng đường, tìm giá trị lớn nhất của hàm bậc hai trên một đoạn), Bài 2"
%% trang_thai: chinh-thuc
%% ghi_chu: "Cùng ý tưởng với L12-B05-0002 (tốc độ thay đổi lớn nhất của mô hình bậc ba, tức giá trị lớn nhất của đạo hàm), nhập vào dạng 12.5.2 đã duyệt."
\begin{ex}
Xét mô hình hóa, một chất điểm chuyển động trên trục số với vị trí cho bởi $s(t)=-t^3+5t^2+4t+6$; trong đó $s$ tính bằng mét, $t$ tính bằng giây $(t\ge0)$. Trong $10$ giây đầu của chuyển động, vận tốc của chất điểm đạt giá trị lớn nhất tại thời điểm nào?
\choice{$t=3$}{\True $t=\dfrac53$}{$t=5$}{$t=8$}
\loigiai{Vận tốc $v(t)=s'(t)=-3t^2+10t+4$ với $t\in[0;10]$. $v'(t)=-6t+10=0\Leftrightarrow t=\dfrac53$. $v(0)=4$, $v\left(\dfrac53\right)=\dfrac{37}3$, $v(10)=-196$. Vậy vận tốc lớn nhất tại $t=\dfrac53$.}
\end{ex}
